\chapter*{EJECUCION DE PARTIDAS Y ANALISIS DE COSTOS UNITARIOS}
\addcontentsline{toc}{chapter}{EJECUCION DE PARTIDAS Y ANALISIS DE COSTOS UNITARIOS}

% Reinicio de contadores para numeracion coherente con el capitulo 4.
\setcounter{section}{0}
\renewcommand{\thesection}{4.\arabic{section}.}
\renewcommand{\thesubsection}{4.\arabic{section}.\arabic{subsection}.}
\renewcommand{\theHsection}{partidas.\arabic{section}}
\renewcommand{\theHsubsection}{partidas.\arabic{section}.\arabic{subsection}}

\section{ALCANCE Y FUENTES DEL ANALISIS}

Los analisis de costos unitarios (ACU) de este capitulo son una base academica para
explicar la ejecucion de las partidas. No constituyen una cotizacion de obra ni una
orden de compra. Los metrados proceden de
\texttt{presupuesto/bom\_final\_aquiles.json} y de los datos electricos canonicos de
los pisos 1 y 2. Los precios de materiales se tomaron del presupuesto fechado el 8 de
junio de 2026; cada recurso se identifica como verificado (V) o estimado (E).

Para mano de obra se adopta, solo con fines academicos, un jornal de S/ 80.00 para
tecnico electricista y S/ 60.00 para ayudante. Las herramientas menores equivalen al
5\% de la mano de obra directa. Los costos no incluyen gastos generales, utilidad,
IGV, transporte extraordinario ni trabajos civiles ajenos a la partida.

\begin{center}
\small
\begin{tabularx}{\textwidth}{p{0.24\textwidth} L}
\toprule
\textbf{Clave} & \textbf{Significado} \\
\midrule
V & Precio comercial con evidencia o enlace registrado el 8 de junio de 2026. \\
E & Precio o rendimiento academico estimado; debe cotizarse antes de comprar o contratar. \\
\bottomrule
\end{tabularx}
\end{center}

\section{METRADO BASE CONCILIADO}

\begin{center}
\small
\begin{tabularx}{\textwidth}{L c r L}
\toprule
\textbf{Partida o recurso} & \textbf{Unidad} & \textbf{Cantidad} & \textbf{Observacion} \\
\midrule
Tuberia PVC SAP de circuitos C1 a C8 & m & 152 & Suma de 16, 20, 11, 24, 28, 22, 16 y 15 m. \\
Conductores de circuitos C1 a C8 & m & 306 & 262 m de 2.5 mm\(^2\) y 44 m de C6. \\
Salidas de alumbrado & pto & 15 & 4 en C1 y 11 en C4. \\
Tomacorrientes generales & pto & 19 & 7 en C2 y 12 en C5. \\
Salidas especiales & pto & 6 & C3: 1; C6: 3; C7: 1; C8: 1. \\
Alimentador principal & m & 40 & Metrado de conductores de 16 mm\(^2\), denominacion monofasica. \\
Canalizacion del alimentador principal & m & 13 & PVC SAP de 40 mm, preliminar. \\
Alimentador TG--T2 & m & 75 & Tres conductores de 10 mm\(^2\), recorrido estimado de 25 m. \\
Canalizacion TG--T2 & m & 25 & PVC SAP de 32 mm, preliminar. \\
Tableros TG y T2 & und & 2 & Capacidad modular pendiente de cierre con unifilar. \\
Sistema de puesta a tierra & cj & 1 & Incluye electrodo, registro, conexion y medicion. \\
\bottomrule
\end{tabularx}
\end{center}

La version rapida revisada indicaba 190 m de canalizacion y cinco salidas especiales.
Esas cantidades no se reproducen porque no coinciden con el BOM: la suma conciliada es
152 m para C1 a C8 y seis salidas especiales. El tubo de C6 se especifica en 32 mm para
el escenario de conductor de 10 mm\(^2\); este cambio debe reflejarse en la proxima
regeneracion del presupuesto.

\section{SECUENCIA GENERAL DE EJECUCION}

\begin{enumerate}
  \item Revisar planos, cuadro de cargas, unifilar, metrados y fichas tecnicas. Cerrar
  las cargas de cocina, lavadora y bomba antes de comprar protecciones o tomas.
  \item Replantear tableros, cajas, rutas y alturas; coordinar interferencias con
  arquitectura, estructura y redes sanitarias.
  \item Ejecutar picado, cajas y canalizaciones. Inspeccionar y probar con guia antes
  del resane o vaciado.
  \item Tender, identificar y conectar conductores con la obra civil seca.
  \item Instalar tableros, protecciones y sistema de puesta a tierra conforme al
  unifilar aprobado.
  \item Ejecutar pruebas, corregir observaciones, actualizar planos conforme a obra y
  emitir el protocolo de puesta en servicio.
\end{enumerate}

\section{ACU-01: CANALIZACION PVC SAP 20 MM EMPOTRADA}

\textbf{Unidad:} m. \textbf{Rendimiento:} 16 m/dia por cuadrilla.

\begin{center}
\small
\begin{tabularx}{\textwidth}{p{0.16\textwidth} L r r r c}
\toprule
\textbf{Tipo} & \textbf{Recurso} & \textbf{Cantidad} & \textbf{P.U. S/} & \textbf{Parcial S/} & \textbf{Base} \\
\midrule
Material & Tubo PVC SAP electrico 20 mm & 1.05 m & 1.97 & 2.07 & V \\
Material & Union PVC SAP 20 mm & 0.35 und & 1.10 & 0.39 & V \\
Material & Curva PVC SAP 20 mm & 0.10 und & 2.20 & 0.22 & V \\
Material & Adhesivo y consumibles & 1 glb & 0.15 & 0.15 & E \\
M.O. & Tecnico electricista & 0.0625 jor & 80.00 & 5.00 & E \\
M.O. & Ayudante & 0.0625 jor & 60.00 & 3.75 & E \\
Equipo & Herramientas menores, 5\% M.O. & 1 glb & 0.44 & 0.44 & E \\
\midrule
\multicolumn{4}{r}{\textbf{Costo unitario directo}} & \textbf{S/ 12.02} & \\
\bottomrule
\end{tabularx}
\end{center}

Incluye replanteo, picado, instalacion, fijacion, resane basico, limpieza y prueba con
guia. No incluye cajas, conductores ni acabados especiales.

\section{ACU-02: CABLEADO 2X2.5 MM\texorpdfstring{$^2$}{2} + PE 2.5 MM\texorpdfstring{$^2$}{2}}

\textbf{Unidad:} m de recorrido. \textbf{Rendimiento:} 40 m/dia por cuadrilla.

\begin{center}
\small
\begin{tabularx}{\textwidth}{p{0.16\textwidth} L r r r c}
\toprule
\textbf{Tipo} & \textbf{Recurso} & \textbf{Cantidad} & \textbf{P.U. S/} & \textbf{Parcial S/} & \textbf{Base} \\
\midrule
Material & Conductor THW 2.5 mm\(^2\), tres conductores y 5\% de reserva & 3.15 m & 2.10 & 6.62 & V \\
Material & Identificacion y terminales ordinarios & 1 glb & 0.15 & 0.15 & E \\
M.O. & Tecnico electricista & 0.025 jor & 80.00 & 2.00 & E \\
M.O. & Ayudante & 0.025 jor & 60.00 & 1.50 & E \\
Equipo & Herramientas menores, 5\% M.O. & 1 glb & 0.18 & 0.18 & E \\
\midrule
\multicolumn{4}{r}{\textbf{Costo unitario directo}} & \textbf{S/ 10.45} & \\
\bottomrule
\end{tabularx}
\end{center}

El ACU corresponde a un recorrido con fase, neutro y PE. Los retornos de alumbrado,
alimentadores y circuitos de seccion diferente se analizan por separado.

\section{ACU-03: SALIDA DE ALUMBRADO SIMPLE}

\textbf{Unidad:} pto. \textbf{Rendimiento:} 6 puntos/dia por cuadrilla.

\begin{center}
\small
\begin{tabularx}{\textwidth}{p{0.16\textwidth} L r r r c}
\toprule
\textbf{Tipo} & \textbf{Recurso} & \textbf{Cantidad} & \textbf{P.U. S/} & \textbf{Parcial S/} & \textbf{Base} \\
\midrule
Material & Caja octogonal & 1 und & 1.90 & 1.90 & V \\
Material & Caja rectangular & 1 und & 1.90 & 1.90 & V \\
Material & Interruptor simple y placa & 1 und & 5.90 & 5.90 & E \\
Material & Luminaria LED plafon 18 W & 1 und & 20.00 & 20.00 & E \\
Material & Conectores y accesorios & 1 glb & 2.00 & 2.00 & E \\
M.O. & Tecnico electricista & 0.15 jor & 80.00 & 12.00 & E \\
M.O. & Ayudante & 0.15 jor & 60.00 & 9.00 & E \\
Equipo & Herramientas menores, 5\% M.O. & 1 glb & 1.05 & 1.05 & E \\
\midrule
\multicolumn{4}{r}{\textbf{Costo unitario directo}} & \textbf{S/ 53.75} & \\
\bottomrule
\end{tabularx}
\end{center}

No incluye tuberia ni conductores, ya valorados en ACU-01 y ACU-02. Para puntos
conmutados se sustituira el interruptor y se agregaran los conductores de control.

\section{ACU-04: TOMACORRIENTE DOBLE 2P+T 16 A}

\textbf{Unidad:} pto. \textbf{Rendimiento:} 8 puntos/dia por cuadrilla.

\begin{center}
\small
\begin{tabularx}{\textwidth}{p{0.16\textwidth} L r r r c}
\toprule
\textbf{Tipo} & \textbf{Recurso} & \textbf{Cantidad} & \textbf{P.U. S/} & \textbf{Parcial S/} & \textbf{Base} \\
\midrule
Material & Tomacorriente doble 2P+T 16 A y placa & 1 und & 19.49 & 19.49 & V \\
Material & Caja rectangular & 1 und & 1.90 & 1.90 & V \\
Material & Terminales y accesorios & 1 glb & 1.50 & 1.50 & E \\
M.O. & Tecnico electricista & 0.10 jor & 80.00 & 8.00 & E \\
M.O. & Ayudante & 0.10 jor & 60.00 & 6.00 & E \\
Equipo & Herramientas menores, 5\% M.O. & 1 glb & 0.70 & 0.70 & E \\
\midrule
\multicolumn{4}{r}{\textbf{Costo unitario directo}} & \textbf{S/ 37.59} & \\
\bottomrule
\end{tabularx}
\end{center}

No incluye tuberia ni conductores. La valorizacion exige prueba de polaridad y
continuidad del conductor PE.

\section{ACU-05: SISTEMA DE PUESTA A TIERRA VERTICAL}

\textbf{Unidad:} cj. \textbf{Rendimiento:} 1 conjunto/dia por cuadrilla, sin considerar
condiciones extraordinarias de suelo.

\begin{center}
\small
\begin{tabularx}{\textwidth}{p{0.16\textwidth} L r r r c}
\toprule
\textbf{Tipo} & \textbf{Recurso} & \textbf{Cantidad} & \textbf{P.U. S/} & \textbf{Parcial S/} & \textbf{Base} \\
\midrule
Material & Varilla cobreada 5/8 pulg x 2.40 m & 1 und & 25.00 & 25.00 & E \\
Material & Conductor de cobre 10 mm\(^2\) & 10 m & 3.50 & 35.00 & E \\
Material & Caja de registro con tapa & 1 und & 35.00 & 35.00 & E \\
Material & Conector certificado para varilla & 1 und & 18.00 & 18.00 & V \\
Material & Tratamiento y consumibles & 1 glb & 60.00 & 60.00 & E \\
M.O. & Tecnico electricista & 1.00 jor & 80.00 & 80.00 & E \\
M.O. & Ayudante & 1.00 jor & 60.00 & 60.00 & E \\
Equipo & Medicion con telurimetro & 1 glb & 80.00 & 80.00 & E \\
Equipo & Herramientas menores, 5\% M.O. & 1 glb & 7.00 & 7.00 & E \\
\midrule
\multicolumn{4}{r}{\textbf{Costo unitario directo}} & \textbf{S/ 400.00} & \\
\bottomrule
\end{tabularx}
\end{center}

La cantidad de conductor se ajustara a la distancia real entre electrodo y tablero. La
partida solo se acepta con protocolo de medicion y resistencia no mayor de 25
\(\Omega\).

\section{ACU-06: PRUEBAS Y PUESTA EN SERVICIO}

\textbf{Unidad:} glb. \textbf{Rendimiento:} 1 instalacion/dia.

\begin{center}
\small
\begin{tabularx}{\textwidth}{p{0.16\textwidth} L r r r c}
\toprule
\textbf{Tipo} & \textbf{Recurso} & \textbf{Cantidad} & \textbf{P.U. S/} & \textbf{Parcial S/} & \textbf{Base} \\
\midrule
Equipo & Uso de megometro, multimetro y telurimetro & 1 glb & 80.00 & 80.00 & E \\
Material & Etiquetas, formatos y consumibles & 1 glb & 15.00 & 15.00 & E \\
M.O. & Tecnico electricista & 0.50 jor & 80.00 & 40.00 & E \\
M.O. & Ayudante & 0.50 jor & 60.00 & 30.00 & E \\
Equipo & Herramientas menores, 5\% M.O. & 1 glb & 3.50 & 3.50 & E \\
\midrule
\multicolumn{4}{r}{\textbf{Costo unitario directo}} & \textbf{S/ 168.50} & \\
\bottomrule
\end{tabularx}
\end{center}

\section{CRITERIO DE VALORIZACION}

Los ACU anteriores no deben sumarse directamente usando cantidades de ``puntos
completos'' que vuelvan a incluir tuberia y conductores. La valorizacion se preparara
con partidas no superpuestas: canalizacion por metro, cableado por metro de recorrido,
dispositivos por punto, tableros por unidad, puesta a tierra por conjunto y pruebas en
forma global.

El costo de tableros y salidas especiales queda pendiente de cierre porque el proyecto
aun debe confirmar las cargas de cocina, lavadora y bomba, la agrupacion de
diferenciales, la corriente nominal aguas arriba y la capacidad modular del gabinete.
Publicar un costo cerrado con los actuales tableros de 8 polos daria una precision
aparente y tecnicamente incorrecta.

\section{CONTROL DE CALIDAD Y CONFORMIDAD}

\begin{center}
\small
\begin{tabularx}{\textwidth}{p{0.25\textwidth} L}
\toprule
\textbf{Control} & \textbf{Criterio de aceptacion} \\
\midrule
Canalizaciones & Sin obstrucciones, aplastamientos ni afectacion estructural; cajas accesibles y alineadas. \\
Cableado & Aislamiento integro, colores identificados, empalmes solo en cajas y continuidad de PE. \\
Tomacorrientes & Polaridad correcta, tierra operativa, placa firme y proteccion diferencial asociada. \\
Tableros & Unifilar aprobado, reserva modular, rotulado, neutro y PE separados, diferenciales coordinados. \\
Puesta a tierra & Conexion accesible y resistencia medida no mayor de 25 \(\Omega\). \\
Pruebas & Protocolos completos, instrumentos identificados y observaciones cerradas. \\
\bottomrule
\end{tabularx}
\end{center}
